% TODO make hw stylesheet
\documentclass[10pt]{article}
\usepackage[letterpaper,top=1in]{geometry}
\usepackage[dvipsnames,svgnames]{xcolor}
\usepackage[shortlabels]{enumitem}
\usepackage[framemethod=TikZ]{mdframed}
\usepackage{amsmath,amssymb,amsthm}
\usepackage[colorlinks]{hyperref}
\usepackage{multicol}
\usepackage{tabularx}

\hypersetup{
	colorlinks=true,
	linkcolor=blue,
	filecolor=magenta,
	urlcolor=magenta,
	pdftitle={MAT 324 HW}
}

\usepackage{mathtools}
\usepackage{thmtools}
\usepackage[textsize=footnotesize]{todonotes}
\usepackage{titling}

\reversemarginpar
\mdfdefinestyle{mdbluebox}{roundcorner=10pt,innerbottommargin=9pt,
	linecolor=blue,backgroundcolor=TealBlue!5,}
\declaretheoremstyle[headfont=\sffamily\bfseries\color{MidnightBlue},
mdframed={style=mdbluebox},]{thmbluebox}

\mdfdefinestyle{mdbloobox}{frametitlefont=\bfseries,innerbottommargin=8pt,
	nobreak=true,backgroundcolor=TealBlue!5,linecolor=blue,}
\declaretheoremstyle[headfont=\bfseries\color{MidnightBlue},
mdframed={style=mdbloobox},headpunct={\\[3pt]},postheadspace=0pt,]{thmbloobox}

\mdfdefinestyle{mdredbox}{frametitlefont=\bfseries,innerbottommargin=8pt,
	nobreak=true,backgroundcolor=Salmon!5,linecolor=RawSienna,}
\declaretheoremstyle[headfont=\bfseries\color{RawSienna},
mdframed={style=mdredbox},headpunct={\\[3pt]},postheadspace=0pt,]{thmredbox}

\mdfdefinestyle{mdgreenbox}{linecolor=ForestGreen,backgroundcolor=ForestGreen!5,
	linewidth=2pt,rightline=false,leftline=true,topline=false,bottomline=false,}
\declaretheoremstyle[headfont=\bfseries\sffamily\color{ForestGreen!70!black},
mdframed={style=mdgreenbox},headpunct={ --- },]{thmgreenbox}

\mdfdefinestyle{mdorangebox}{linecolor=RedOrange,backgroundcolor=RedOrange!5,
	linewidth=2pt,rightline=false,leftline=true,topline=false,bottomline=false,}
\declaretheoremstyle[headfont=\bfseries\sffamily\color{RedOrange!70!black},
mdframed={style=mdorangebox},headpunct={ --- },]{thmorangebox}

\mdfdefinestyle{mdblackbox}{linecolor=black,backgroundcolor=RedViolet!5!gray!5,
	linewidth=3pt,rightline=false,leftline=true,topline=false,bottomline=false,}
\declaretheoremstyle[mdframed={style=mdblackbox}]{thmblackbox}

\declaretheoremstyle[spaceabove=6pt,spacebelow=6pt]{basehead}

\declaretheorem[style=basehead,name=Problem]{problem}
\declaretheorem[style=thmredbox,name=Problem,numbered=no]{problem*}
\declaretheorem[style=thmbloobox,name=Setup]{setup} % for config geo
\declaretheorem[style=thmredbox,name=Example,sibling=problem]{example}
\declaretheorem[style=thmredbox,name=$\bigstar$ Problem,sibling=problem]{niceprob}
\declaretheorem[style=thmbluebox,name=Theorem,sibling=problem]{theorem}
\declaretheorem[style=thmbluebox,name=Corollary,sibling=theorem]{corollary}
\declaretheorem[style=thmbluebox,name=Proposition,sibling=theorem]{prop}
\declaretheorem[style=thmbluebox,name=Lemma,numberwithin=problem]{lemma}
\declaretheorem[style=thmbluebox,name=Theorem,numbered=no]{theorem*}
\declaretheorem[style=thmbluebox,name=Lemma,numbered=no]{lemma*}
\declaretheorem[style=thmgreenbox,name=Claim,numberwithin=problem]{claim}
\declaretheorem[style=thmgreenbox,name=Claim,numbered=no]{claim*}
\declaretheorem[style=thmblackbox,name=Remark,numberwithin=problem]{remark}
\declaretheorem[style=thmblackbox,name=Remark,numbered=no]{remark*}
\declaretheorem[style=thmgreenbox,name=Definition,sibling=theorem]{definition}
\declaretheorem[style=thmgreenbox,name=Definition,numbered=no]{definition*}

\providecommand{\ol}{\overline}
\providecommand{\ceil}[1]{\left\lceil #1 \right\rceil}
\providecommand{\floor}[1]{\left\lfloor #1 \right\rfloor}
\newcommand{\dd}{\mathrm{d}}
\providecommand{\ul}{\underline}
\providecommand{\eps}{\varepsilon}
\providecommand{\half}{\frac{1}{2}}
\providecommand{\CC}{\mathbb C}
\providecommand{\FF}{\mathbb F}
\providecommand{\NN}{\mathbb N}
\providecommand{\QQ}{\mathbb Q}
\providecommand{\RR}{\mathbb R}
\providecommand{\ZZ}{\mathbb Z}
\providecommand{\dg}{^\circ}
\providecommand{\ii}{\item}
\providecommand{\setdiff}{\smallsetminus}
\providecommand{\alert}[1]{{\sffamily\textbf{\textcolor{blue}{#1}}}}
\providecommand{\inv}{^{-1}}
\providecommand{\mb}[1]{\mathbf{#1}}

\newenvironment{soln}{\begin{proof}[Solution]}{\end{proof}}

\providecommand{\pow}{\operatorname{Pow}}
\providecommand{\vocab}[1]{{\textbf{\textcolor{ForestGreen}{#1}}}}
\providecommand{\lcm}{\operatorname{lcm}}
\providecommand{\abs}[1]{\left|#1\right|}

\usepackage{mathrsfs}

% place title at top right
\setlength{\droptitle}{-8ex}
\pretitle{\begin{flushright}\Large\bfseries}
\posttitle{\par\end{flushright}}
\preauthor{\begin{flushright}}
\postauthor{\end{flushright}}
\predate{\begin{flushright}}
\postdate{\end{flushright}}

\title{MAT 324 HW03}
\author{\large Aditya Pahuja}
\date{\large Due 09/18}
\begin{document}
\maketitle

\begin{problem}
    Let $(X,d_X)$ and $(Y,d_Y)$ be metric spaces, $A\subseteq X$,
    and $f\colon X\to Y$ a map.
    \begin{enumerate}[(a)]
        \ii Suppose $f$ is continuous at every $x\in A$.
	Show that $f|_A$ is continuous.
	\ii Given an example showing that the converse need not hold,
	i.e. $f|_A$ may be continuous despite $f$ being discontinuous
	at some $x\in A$.
    \end{enumerate}
\end{problem}

\begin{soln}
    (a) If $f$ is continuous on $X$ then, for every open $V\subseteq Y$,
    $f\inv(V)$ is open in $X$, so
    $f|_A\inv(V) = A\cap f\inv(V)$ is open in $A$;
    this is what it means for $f|_A$ to be continuous.\\

    (b) Take $f\colon\RR\to\RR$ defined as
    \begin{equation*}
        f(x) = \begin{cases}
	    0 & \text{if }x<0\\
	    1 & \text{if }x\ge1
        \end{cases}
    \end{equation*}
    and let $A = [0,\infty)$. Then, $f|_A$ is constant, hence continuous,
    but $f$ is not continuous at $0\in A$ because $f(0)\in(\half,\infty)$ but
    $f\inv((\half,\infty)) = [0,\infty)$ does not contain a neighborhood of $0$.
\end{soln}

\begin{problem}
    Suppose $A\subseteq\RR$ is closed and $f\colon A\to\RR$ is a continuous function.
    Thus, $\RR\setminus A$ is a union of at most countably many disjoint intervals
    $(a_k,b_k)$ with $a_k,b_k\in A\sqcup\{\pm\infty\}$ distinct for $k\in\NN$.
    Show that the function $g\colon\RR\to\RR$ defined by
    \begin{equation*}
	g(x) = \begin{cases}
	    f(x) & \text{if }x\in A\\
	    f(b_k) & \text{if }x\in(a_k,b_k)\text{ for some }k
	    \text{ such that }a_k=-\infty\\
	    f(a_k) & \text{if }x\in(a_k,b_k)\text{ for some }k
	    \text{ such that }b_k=\infty\\
	    tf(a_k) + (1-t)f(b_k) & \text{if }x = ta_k + (1-t)b_k
	    \text{ for some }t\in[0,1]\text{ and }k
	\end{cases}
    \end{equation*}
    is continuous.
\end{problem}

\begin{soln}
    If $x$ is not a boundary point of $A$ (i.e. it either lies outside $A$
    or in the interior of $A$) then $g$ is continuous at $x$ because
    $g$ is continuous on the open sets $\operatorname{Int}(A)$ and $\RR\setminus A$.

    Now suppose $x\in \partial A$ is on the boundary. We consider two cases.
    First, if there is no sequence of $b_k$s converging to $x$,
    then there is a neighborhood $B = (x-\delta,x+\delta)$
    of $x$ not containing any boundary points besides possibly $x$.
    By construction $g$ must be continuous on $(x-\delta, x]$ and $[x,x+\delta)$,
    since $g\equiv f$ on one of these intervals and $g$ is described by
    a linear polynomial on the other (because one of the intervals is contained in $A$
    and the other isn't); the pasting lemma then gives continuity
    on the neighborhood $B$, so $g$ is continuous at $x$.

    The alternative is that some sequence of $b_k$s converges to $x$;
    call this sequence $(\widehat b_n)_{n\in\NN}$
    and let $\widehat a_n$ be the left endpoint of the interval
    in $\RR\setminus A$ that has $\widehat b_n$ as its right endpoint.
    Assume that the $\widehat b_n$s are all less than $x$; we will show that
    \begin{equation*}
	\lim_{\substack{x'\to x^-}} g(x') = g(x)
    \end{equation*}
    The sequence $\widehat a_n$ also converges to $x$ because
    \begin{equation*}
	\sup_{\widehat b_k \le \widehat a_n} \widehat b_k
	\le \widehat a_n \le \widehat b_n.
    \end{equation*}
    Now we have for any $\delta>0$ that if $x'\in(x-\delta,x]\setminus A$, then
    \begin{equation*}
	\abs{g(x)-g(x')}\le
	\max\left(\sup_{\widehat a_k\in (x-\delta,x]} \abs{g(x) - g(\widehat a_k)},
	\sup_{\widehat b_k\in (x-\delta,x]} \abs{g(x)-g(\widehat b_k)}\right)
    \end{equation*}
    because for each $[\widehat a_k,\widehat b_k]\subseteq (x-\delta,x]$,
    $g$ reaches its maximum and minimum value at the endpoints $a_k$, $b_k$.
    Since the $\widehat a_k$ and $\widehat b_k$ converge to $x$,
    given $\eps>0$ we can pick $\delta$ so that the upper bound is
    less than $\eps$, which implies that $\abs{g(x)-g(x')} < \eps$.
    By continuity of $f$ at $x$, we can also shrink $\delta$ to ensure that
    if $x'\in (x-\delta,x]\cap A$ then $\abs{g(x)-g(x')} < \eps$
    still holds, hence the bound holds for all $x'\in (x-\delta,x]$. Thus
    \begin{equation*}
	\lim_{x'\to x^-} g(x') = g(x).
    \end{equation*}

    To get the limit from above, we can find a sequence of $a_k$s,
    say $\widetilde a_n$, (and the corresponding $\widetilde b_n$) whose terms are
    all at least as large as $x$ that converges to $x$.
    Then, consider $h(x) = g(-x)$, which is now defined
    as a extension of $x\mapsto f(-x)$ by linear polynomials on intervals
    $(-b_k,-a_k)$. The $-\widetilde a_n$ and $-\widetilde b_n$
    converge to $-x$ from below and $-x$ is in the closed set
    $-A = \{-a : a\in A\}$ on which $x\mapsto f(-x)$
    is defined, so the preceding paragraph applies to $h$
    and we may use continuity from below
    for $h$ at $-x$ to extract continuity from above for $g$ at $x$.
    Thus $g$ is continuous at all $x\in A$, completing the proof.
\end{soln}

\begin{problem}
    Let $A\subseteq\RR$. Show that the function $f\colon\RR\to\RR$ given by
    \begin{equation*}
        f(x) = \begin{cases}
	    x & \text{if }x\in A\\
	    -x & \text{if }x\notin A
        \end{cases}
    \end{equation*}
    is Lebesgue measurable if and only if $A$ is a Lebesgue measurable set.
\end{problem}

\begin{soln}
    Since $\{(a,\infty) : a\in\RR\}$ generates the Borel $\sigma$-algebra on $\RR$,
    $f$ is Lebesgue measurable if and only if $f\inv((a,\infty))$ is Lebesgue
    for all $a\in\RR$. To show that this is equivalent to $A$ being Lebesgue,
    we write
    \begin{equation*}
        f\inv((a,\infty)) = (A\cap(a,\infty)) \cup (A^c\cap(-\infty,-a)).
    \end{equation*}
    If $A$ is Lebesgue, then this set is of course Lebesgue for all real $a$
    because Lebesgue sets form a $\sigma$-algebra.
    Conversely, if $f\inv((a,\infty))$ is Lebesgue for all real $a$, then
    \begin{equation*}
	A\cap(0,\infty) = (0,\infty)\cap f\inv((0,\infty)),
	\quad
	A\cap(-\infty,0) = (-\infty,0)\cap f\inv((0,\infty)),
    \end{equation*}
    are both Lebesgue, implying
    \begin{equation*}
	A\setminus\{0\}
	= (A\cap(0,\infty))\cup ((-\infty,0)\setminus(A^c\cap(-\infty,0))
    \end{equation*}
    is Lebesgue. If $A\setminus\{0\} = A$ then $A$ is Lebesgue,
    otherwise $(A\setminus\{0\}) \cup\{0\} = A$ so $A$ is still Lebesgue.
    Thus $A$ is Lebesgue if and only if the pre-image
    of every open ray under $f$ is Lebesgue.
\end{soln}

\pagebreak

\begin{problem}
    Let $(X,\mathcal S)$ be a measurable space.
    \begin{enumerate}[(a)]
        \ii Suppose $(Y,d)$ is a metric space,
	$g\colon Y\to\RR$ is a continuous function,
	and $f\colon X\to Y$ is a function such that $f\inv(V)\in\mathcal S$
	for every open subset $V\subseteq Y$.
	Show that the function $g\circ f\colon X\to\RR$
	is $\mathcal S$-measurable.
	\ii Suppose $f_1,\dots,f_k\colon X\to\RR$ are
	$\mathcal S$-measurable functions and
	$f = (f_1,\dots, f_k)\colon X\to\RR^k$.
	Show that $f\inv(V)\in\mathcal S$ for every
	open subset $V\subseteq\RR$.
	\ii Show that $f_1,f_2\colon X\to\RR$ are
	$\mathcal S$-measurable functions and $a_1,a_2\in\RR$.
	Show that
	\begin{equation*}
	    a_1f_1 + a_2f_2\colon X\to\RR\text{ and }f_1f_2\colon X\to\RR
	\end{equation*}
	are also $\mathcal S$-measurable functions.
	If in addition $f_2\inv(0) = \varnothing$, show that
	$f_1/f_2\colon X\to\RR$ is an $\mathcal S$-measurable function.
	\ii Show that a function $f\colon X\to\RR$ is Lebesgue measurable
	if and only if the function $f^3$ is Lebesgue measurable.
	\ii Give an example of a function $f\colon\RR\to\RR$
	so that it is not Lebesgue measurable,
	but $f^2\colon\RR\to\RR$ is Lebesgue measurable.
    \end{enumerate}
\end{problem}

\begin{soln}
    (a) It suffices to check that the pre-image of an open set
    under $g\circ f$ is in $\mathcal S$.
    Suppose $U\subseteq\RR$ is open. Then
    $g\inv(U)$ is open in $Y$ by continuity,
    so $(g\circ f)\inv(U) = f\inv(g\inv(U))$ is in $\mathcal S$.\\

    (b) The rational cubes $(a_1,b_1)\times\cdots\times(a_k,b_k)$,
    where $a_j<b_j$ for $j=1,\dots,k$ are rational numbers,
    are a basis for the topology of $\RR^k$.
    That is, for any open $V\subseteq\RR^k$, we can find cubes
    $I_n = (a^{(n)}_1, b^{(n)}_1)\times\cdots\times(a^{(n)}_k,b^{(n)}_k)$,
    where the intervals have endpoints in $\QQ$,
    such that $V = \bigcup_{n=1}^\infty I_n$. Then,
    \begin{equation*}
	f\inv(V) = \bigcup_{n=1}^\infty f\inv(I_n) =
	\bigcup_{n=1}^\infty \bigcap_{j=1}^k f_j\inv((a^{(n)}_j,b^{(n)}_j))
    \end{equation*}
    because the pre-image of $I_n$ is the set of $x\in X$ such that
    $f_j(x)\in(a^{(n)}_j,b^{(n)}_j)$ for each $j=1,\dots,k$.
    In summary, $f\inv(V)$ is the union of countably many
    finite intersections of $\mathcal S$-measurable sets,
    since of course $f_j\inv((a_j^{(n)},b_j^{(n)}))$ is $\mathcal S$-measurable,
    so $f\inv(V)$ is $\mathcal S$-measurable.\\

    (c) Define $f\colon X\to\RR^2$ by $f(x) = (f_1(x), f_2(x))$;
    this is $\mathcal S$-measurable because, by part (b),
    the pre-image of every open subset of $\RR$ under $f$ is in $\mathcal S$,
    which by part (a) implies that $f$ is a $\mathcal S$-measurable function.
    Then, the functions $g_1,g_2\colon\RR^2\to\RR$ and
    $g_3\colon\RR\times(\RR\setminus\{0\})\to\RR$ given by
    \begin{equation*}
        g_1(x,y) = a_1x+a_2y,\quad g_2(x,y) = xy,\quad g_3(x,y) = x/y
    \end{equation*}
    are all continuous on their domains, so by (a)
    the functions $g_1\circ f = a_1f_1 + a_2f_2$
    and $g_2\circ f = f_1f_2$ are $\mathcal S$-measurable,
    and if $f_2\inv(0)=\varnothing$,
    then $g_3\circ f = f_1/f_2$ is also $\mathcal S$-measurable.\\

    (d) If $f\colon X\to\RR$ is measurable then we can compose $f$
    with the continuous map $x\mapsto x^3$ to get that $f^3$ 
    is also measurable by (a). Conversely if $f^3\colon X\to\RR$ is measurable
    then we compose $f^3$ with the continuous map $x\mapsto\sqrt[3]{x}$
    to get that $\sqrt[3]{f^3} = f$ is also measurable by (a).\\

    (e) Let $V\subseteq\RR$ be a non-Lebesgue set.
    The function $f\colon\RR\to\RR$ defined by
    \begin{equation*}
        f(x) = \begin{cases}
	    1 & \text{if }x\in V\\
	    -1 & \text{if }x\notin V.
        \end{cases}
    \end{equation*}
    Then the pre-image of the Borel set $\{1\}$ is $V$,
    so $f$ is not Lebesgue measurable, but $f^2\equiv1$ everywhere
    so obviously it is Lebesgue measurable.
\end{soln}

\pagebreak
\begin{problem}
    Let $(X,\mathcal S, \mu)$ be a measurable space and
    $f_1,f_2,\dots\colon X\to[-\infty,\infty]$ be
    $\mathcal S$-measurable functions converging to
    a function $f\colon X\to[-\infty,\infty]$ pointwise.
    Thus, $f$ is also $\mathcal S$-measurable and the subsets
    $X^\pm = f\inv(\pm\infty)$ are elements of $\mathcal S$.
    \begin{enumerate}[(a)]
        \ii Suppose $\mu(X^+) < \infty$.
	Show that for every $\eps>0$ there exists $X_\eps^+\in\mathcal S$
	so that $X_\eps^+\subseteq X^+$, $\mu(X^+\setminus X_\eps^+)<\eps$,
	and the functions $f_1|_{X_\eps^+}$, $f_2|_{X_\eps^+}$, \dots
	converge uniformly to $\infty$.
	\ii Suppose $\mu(X) < \infty$. Show that for every $\eps > 0$
	there exists $X_\eps\in\mathcal S$ so that $\mu(X\setminus X_\eps)<\eps$
	and the functions $f_1|_{X_\eps}$, $f_2|_{X_\eps}$, \dots
	converge uniformly to $f|_{X_\eps}$.
	\ii Give an example showing that the last conclusion can fail
	if $\mu(X) = \infty$.
    \end{enumerate}
\end{problem}

\begin{soln}
    (a) For $N,t\in\NN$, let $X_{N,t} = X^+\cap
    \bigcap_{n=N}^\infty f_n\inv([t,\infty])$;
    this is the subset of $X^+$ on which
    $f_n(x) \ge t$ for all $n\ge N$. Then,
    \begin{equation*}
	X^+ = \bigcup_{N=1}^\infty X_{N,t}
    \end{equation*}
    because for each $x\in X^+$, given $t\in\NN$, there exists $N\in\NN$
    such that $n\ge N$ implies $f_n(x)\ge t$ ---
    this is what pointwise convergence means.
    Clearly $X_{N,t}\subseteq X_{N+1,t}$ for all $N$ and $t$,
    so monotone convergence of measure gives
    \begin{equation*}
	\mu(X^+) = \lim_{N\to\infty} \mu(X^+\cap X_{N,t})
    \end{equation*}
    for all $t>0$. The limit is finite, so for each positive integer $t$
    there is an $N_t$ such that $\mu(X^+)$ is within $\eps/2^t$ of
    $\mu(X_{N_t,t})$, hence $\mu(X^+\setminus X_{N_t,t}) < \eps/2^t$. Let
    \begin{equation*}
	X^+_\eps = \bigcap_{t=1}^\infty X_{N_t,t}.
    \end{equation*}
    For each $t\in\NN$, $X^+_\eps\subseteq X_{N_t,t}$, which means that
    there is an $N=N_t\in\NN$ depending only on $t$ such that
    $f_n(x)\ge t$ for all $n\ge N$, i.e. convergence of $f_1$, $f_2$, \dots
    on $X^+_\eps$ is uniform. Furthermore,
    \begin{equation*}
	\mu(X^+\setminus X^+_\eps) \le
	\mu\left(\bigcup_{t=1}^\infty X^+\setminus X_{N_t,t}\right) \le
	\sum_{t=1}^\infty \mu(X^+\setminus X_{N_t,t}) < \sum_{t=1}^\infty \eps/2^t
	= \eps,
    \end{equation*}
    so $X^+_\eps$ has the desired properties.
    \\

    (b) Let $X^0 = f\inv(\RR)\in\mathcal S$
    be the set of points on which $f$ is finite.
    By Egorov's theorem for sequences of functions
    that converge to a function mapping into the reals,
    there exists $X^0_{\eps/3}\subseteq X^0$ in $\mathcal S$ such that
    $\mu(X^0\setminus X^0_{\eps/3})<\eps/3$ and
    $f_1|_{X^0_{\eps/3}}$, $f_2|_{X^0_{\eps/3}}$, \dots converges
    uniformly to $f|_{X^0_{\eps/3}}$.

    By part (a), there exists $X^+_{\eps/3}\subseteq X^+$ in $\mathcal S$
    such that $\mu(X^+\setminus X^+_{\eps/3}) < \eps/3$ and the sequence
    $f_1|_{X^+_{\eps/3}}$, $f_2|_{X^+_{\eps/3}}$, \dots converges
    uniformly to $f|_{X^+_{\eps/3}}$.

    Similarly, applying (a) to the sequence $-f_1|_{X^-}$, $-f_2|_{X^-}$, \dots
    converging to $-f|_{X^-}\equiv\infty$ yields the existence of a set
    $X^-_{\eps/3}\subseteq X^-$ in $\mathcal S$
    such that $\mu(X^-\setminus X^-_{\eps/3}) < \eps/3$ and the sequence
    $-f_1|_{X^-_{\eps/3}}$, $-f_2|_{X^-_{\eps/3}}$, \dots converges
    uniformly to $-f|_{X^-_{\eps/3}}$, or equivalently
    $f_1|_{X^-_{\eps/3}}$, $f_2|_{X^-_{\eps/3}}$, \dots converges
    uniformly to $f|_{X^-_{\eps/3}}$.

    Then, take $X_\eps = X^0_{\eps/3}\cup X^+_{\eps/3}\cup X^-_{\eps/3}$, so that
    \begin{equation*}
        \mu(X\setminus X_\eps) \le
	\mu(X\setminus X^0_{\eps/3}) +
	\mu(X\setminus X^+_{\eps/3}) +
	\mu(X\setminus X^-_{\eps/3}) <
	\frac\eps3 + \frac\eps3 + \frac\eps3 = \eps
    \end{equation*}
    and we have uniform convergence on each of $X_\eps$
    due to having uniform convergence on the component sets
    $X^0_{\eps/3}$, $X^0_{\eps/3}$, $X^0_{\eps/3}$.
    \\

    (c) Let $f_1,f_2,\dots\colon\RR\to\RR$ be a sequence of functions
    where $f_n$ is the indicator function $\mathbf 1_{[n,n+1]}$.
    Then, $f_n$ converges pointwise to the zero function $f$,
    but, if $\eps<1$, then any $X_\eps$ such that $\abs{\RR\setminus X_\eps}<\eps$
    intersects each $[n,n+1]$ interval (otherwise $\RR\setminus X_\eps$
    contains an interval of length $1$, hence has outer measure bigger than $\eps$),
    so for any $n\in\NN$ we have
    \begin{equation*}
	\sup_{x\in X_\eps} |f(x) - f_n(x)| = 1,
    \end{equation*}
    meaning the convergence to $f$ is not uniform.
\end{soln}

\begin{problem}
    Let $\QQ = \{q_1,q_2,\dots\}$, $I_k = (q_k-2^{-k},q_k+2^{-k})$
    for $k\in\NN$, $U = \bigcup_{k=1}^\infty I_k$, and $A\subseteq\RR$.
    \begin{enumerate}[(a)]
	\ii Suppose $\abs{\RR\setminus A} = 0$. Show that
	$A\cap I_k\ne\varnothing$ for every $k\in\NN$ and
	$A\cap U$ is dense in $A$.
	\ii Suppose $\mathbf 1_U|_A$ is continuous. Show that
	$\abs{\RR\setminus A} > 0$.
    \end{enumerate}
\end{problem}

\begin{soln}
    (a) We show the contrapositive:
    if $A\cap I_k = \varnothing$ for some $k\in\NN$,
    then $I_k\subseteq\RR\setminus A$,
    so $\abs{\RR\setminus A}\ge \abs{I_k} = 2\cdot 2^{-k}$ is nonzero.

    Suppose that $V$ is a nonempty open subset of $A$, so
    $V = V'\cap A$ for some nonempty open subset $V'$ of $\RR$.
    Then, for each $x\in V'$, there is an $\eps>0$ such that
    the open ball $B_\eps(x)$ centered at $x$ is a subset of $V'$
    Choose $N\in\NN$ such that $2^{-N} < \eps/2$, so that
    $B_{2^{-N}}(x)\subseteq B_\eps(x)$.
    Since $B_{2^{-N}}(x)$ contains infinitely many rational numbers,
    there exists an $n\ge N$ such that $q_n\in B_{2^{-N}}(x)$.
    In particular, any point $y\in B_{2^{-n}}(q_n) = I_n$ satisfies
    \begin{equation*}
	\abs{x - y} \le \abs{x-q_n} + \abs{q_n-y} <
	2^{-N} + 2^{-n} \le 2\cdot 2^{-N} < \eps,
    \end{equation*}
    so $V'$ contains $I_n$. Since $A\cap I_n$ is nonempty, $A$ contains a point of
    $U\cap V'\supseteq I_n$, so $U\cap A\cap V' = U\cap A\cap V$ is nonempty,
    which is what it means for $U\cap A$ to be dense in $A$.

    (b) Suppose for the sake of contradiction that
    $\abs{\RR\setminus A} = 0$, so $\abs A$ is infinite
    (else $\abs{A\cup(\RR\setminus A)} = \abs\RR$ would be finite).
    The indicator function of $U$ restricted to $A$
    is the same thing as the function $\mathbf 1_{U\cap A}\colon A\to\RR$.
    Since $\mathbf 1_{U\cap A}$ is constant on $U\cap A$,
    it is uniformly continuous on $U\cap A$, so, since $U\cap A$ is dense in $A$
    by part (a) (and our assumption that $\abs{\RR\setminus A} = 0$),
    $\mathbf 1_{U\cap A}$ uniquely extends to a (uniformly)
    continuous function on $A$; this is an easy result from MAT 320.

    One continuous function $A\to\RR$ whose restriction to $U\cap A$
    is $\mathbf 1_{U\cap A}$ is the constant function $\mathbf 1_A\colon A\to\RR$.
    By uniqueness it must be the only such function,
    implying $\mathbf 1_U|_A$ is the same function as $\mathbf 1_A|_A$,
    which forces $U\cap A = A$, so $A\subseteq U$. However,
    \begin{equation*}
	\abs U \le \sum_{k=1}^\infty \abs{I_k} =
	\sum_{k=1}^\infty 2\cdot 2^{-k} = 2,
    \end{equation*}
    so $\abs A$ is finite because $A$ is contained in a set of finite measure,
    which is a contradiction.
\end{soln}










\end{document}
